\documentclass[11pt,a4paper]{article}
\usepackage{fontspec,amsmath,amssymb,amsthm,microtype}
\setmainfont{lmroman10-regular.otf}[BoldFont=lmroman10-bold.otf,ItalicFont=lmroman10-italic.otf,BoldItalicFont=lmroman10-bolditalic.otf]
\setmonofont{lmmono10-regular.otf}[Scale=MatchLowercase]
\usepackage[margin=25mm]{geometry}
\usepackage[unicode,hidelinks]{hyperref}
\setlength{\emergencystretch}{3em}
\setlength{\parskip}{0.4em}
\setlength{\parindent}{0pt}
\allowdisplaybreaks[2]
\raggedbottom
\title{\texorpdfstring{Actual Gaussian ULA laws and their frozen-\allowbreak{}drift interpolation}{Actual Gaussian ULA laws and their frozen-drift interpolation}}
\author{LangevinLean mathematical exposition}
\date{}
\begin{document}\maketitle
\noindent Source module: \texttt{Unadjusted\allowbreak{}Langevin/\allowbreak{}LSI/\allowbreak{}ULALaw.\allowbreak{}lean}. The accompanying HTML pairs this mathematical exposition with the exact Lean statements and records the scope of the formalization.


Define the frozen-drift law \(\rho_{k,t}=\operatorname{Law}(X_k-t\nabla U(X_k)+\sqrt{2t}Z)\), with independent standard Gaussian \(Z\), and \(\rho_{k+1}=\rho_{k,h}\). The endpoint and probability identities connect the entropy argument to the actual algorithm.


\section{\texorpdfstring{ula\allowbreak{}Entropy}{ulaEntropy}}


Set \(H_k=\operatorname{KL}(\rho_k\mid\nu)_{\mathbb R}\). The definition uses the actual extended KL of the ULA law, followed by \texttt{to\allowbreak{}Real}. Separate theorems prove finite KL wherever this is used as ordinary real entropy; the definition alone would not exclude infinity.



{\raggedright\noindent Lean declarations: \texttt{ula\allowbreak{}Entropy}.\par}

\section{\texorpdfstring{ula\allowbreak{}Entropy zero}{ulaEntropy zero}}


Substitute the definitional identity \(\rho_0=\mu\) into the entropy definition. This gives \(H_0=\operatorname{KL}(\mu\mid\nu)_{\mathbb R}\), the initial term appearing in the final exponential bound.



{\raggedright\noindent Lean declarations: \texttt{ula\allowbreak{}Entropy\_\allowbreak{}zero}.\par}

\section{\texorpdfstring{interpolation\allowbreak{}Law}{interpolationLaw}}


Define \(\rho_t=(T_t)_*(\mu\otimes\gamma)\), where \(T_t(x,z)=x-t\nabla U(x)+\sqrt{2t}z\). The product encodes independence. This is an alias of the actual \texttt{ula\allowbreak{}Law\allowbreak{}Step}; its frozen-drift Gaussian interpolation interpretation applies to \(t\ge0\).



{\raggedright\noindent Lean declarations: \texttt{interpolation\allowbreak{}Law}.\par}

\section{\texorpdfstring{ula\allowbreak{}Laws}{ulaLaws}}


Iterate the actual independent-Gaussian transition at fixed step \(h\): \(\rho_k=P_h^k\mu\). The definition is on measures, while the separate chain-law theorems identify independent iid implementations with these same laws.



{\raggedright\noindent Lean declarations: \texttt{ula\allowbreak{}Laws}.\par}

\section{\texorpdfstring{interpolation\allowbreak{}Law endpoint}{interpolationLaw endpoint}}


At elapsed time \(h\), the interpolation from \(\rho_k\) is \(P_h\rho_k\). The preceding successor identity therefore proves \(\rho_{k,h}=\rho_{k+1}\). This is the law equality used to turn a continuous interpolation bound into the discrete entropy step.



{\raggedright\noindent Lean declarations: \texttt{interpolation\allowbreak{}Law\_\allowbreak{}endpoint}.\par}

\section{\texorpdfstring{ula\allowbreak{}Laws is\allowbreak{}Probability}{ulaLaws isProbability}}


Induct from the initial probability \(\mu\). At the successor rewrite the iterate as one interpolation step at time \(h\), then apply probability preservation. Every algorithm law is therefore a probability before applying KL and LSI identities.



{\raggedright\noindent Lean declarations: \texttt{ula\allowbreak{}Laws\_\allowbreak{}is\allowbreak{}Probability}.\par}

\section{\texorpdfstring{interpolation\allowbreak{}Law zero}{interpolationLaw zero}}


At \(t=0\), both the drift increment and Gaussian increment vanish, so the update map is first-coordinate projection. Integrating out the standard Gaussian probability returns \(\mu\). This endpoint identity needs no potential regularity.



{\raggedright\noindent Lean declarations: \texttt{interpolation\allowbreak{}Law\_\allowbreak{}zero}.\par}

\section{\texorpdfstring{ula\allowbreak{}Laws zero}{ulaLaws zero}}


The zeroth iterate of a function is the identity, hence \(\rho_0=\mu\). Lean uses definitional equality, which makes the initial-entropy rewrite exact rather than an asymptotic endpoint argument.



{\raggedright\noindent Lean declarations: \texttt{ula\allowbreak{}Laws\_\allowbreak{}zero}.\par}

\section{\texorpdfstring{ula\allowbreak{}Laws succ}{ulaLaws succ}}


The successor-iterate identity gives \(\rho_{k+1}=P_h(\rho_k)\). Lean invokes \texttt{Function.\allowbreak{}iterate\_\allowbreak{}succ\_\allowbreak{}apply'}, whose order matches applying one more transition to the current law.



{\raggedright\noindent Lean declarations: \texttt{ula\allowbreak{}Laws\_\allowbreak{}succ}.\par}

\section{\texorpdfstring{interpolation\allowbreak{}Law is\allowbreak{}Probability}{interpolationLaw isProbability}}


Lipschitzness makes the gradient continuous and the joint update measurable. The product \(\mu\otimes\gamma\) is a probability, and its pushforward has the same total mass. Evaluating the mapped measure on the whole space proves normalization of \(\rho_t\).



{\raggedright\noindent Lean declarations: \texttt{interpolation\allowbreak{}Law\_\allowbreak{}is\allowbreak{}Probability}.\par}

\end{document}
